\section{Folds, predictions, and the test--cost theorem}\label{sec:preservation}
This section proves the machinery of the lower bounds in the order of the
overview. We construct the reference structure (\S\ref{sec:structure}) and
define color-determined relations, folds, and predictions
(\S\ref{sec:types}). We then show
that a relational circuit computes the predicted values on every folded
structure whose fold none of its compositions notices (\S\ref{sec:unnoticed}),
and that each composition notices at most two folds
(\S\ref{sec:responses}). With the tests of \S\ref{sec:G}, this already
proves the composition bounds of Theorem~\ref{thm:main}.
Section~\ref{sec:simultaneous} states the test--cost theorem, a counting
form of the same argument that also covers randomized circuits and the
other targets of \S\ref{sec:additional-results};
Section~\ref{sec:matrix} supplies its hypotheses for scalar-valued
computations.

\subsection{The reference structure}\label{sec:structure}
Fix a palette $\Lambda$ of $N\ge2$ colors and a vertex set $V$ of size $n$.
A coloring $\chi:\binom V2\to\Lambda$ assigns a color to every unordered
pair of distinct vertices. We write $\chi(u,v)$ for $\chi(\{u,v\})$.

\Needspace{7\baselineskip}
\begin{lemma}[Two witnesses]\label{lem:palette}
For every integer
\[
 n\ge n_0(N)\coloneqq\left\lceil12N^2\ln(2N)\right\rceil+3,
\]
there is a coloring such that, for all distinct $u,v\in V$ and all
$\mu,\nu\in\Lambda$, at least two distinct vertices
$z\in V\setminus\{u,v\}$ satisfy
$\chi(u,z)=\mu$ and $\chi(z,v)=\nu$.
\end{lemma}
\begin{proof}
Color the edges independently and uniformly. For fixed distinct $u,v,w$
and colors $\mu,\nu$, the probability that no suitable witness occurs
outside $\{u,v,w\}$ is at most $\exp(-(n-3)/N^2)$: the candidate witnesses
use disjoint pairs of edges. A union bound gives failure probability at most
\[
 n^3N^2\exp(-(n-3)/N^2).
\]
At $n=n_0(N)$, the inequality $n_0(N)\le13N^3$ bounds this expression by
$2197/(4096N)<1$. It decreases for $n>3N^2$, proving the assertion at every
larger size. Hence some coloring has, for all distinct $u,v,w$ and all
colors $\mu,\nu$, a witness outside $\{u,v,w\}$; applying this with $w$
equal to one witness for $u,v$ yields a second.
\end{proof}

Fix such a coloring, and choose a vertex $\hub\in V$, called the
hub. The two witnesses of Lemma~\ref{lem:palette} play different
roles. One witness outside the endpoints already makes every computed
relation a union of color classes, possibly together with the diagonal
(Lemma~\ref{lem:tables}). The second ensures that two vertices other
than $\hub$ have, for every pair of colors, a witness that also avoids
$\hub$, so that recoloring the edges at $\hub$ does not affect pairs away
from it (\S\ref{sec:responses}). Lemma~\ref{lem:palette} also implies that
every vertex sees every color, and that every vertex other than $\hub$
sees every color on edges avoiding $\hub$ (take $v=\hub$). For each color
$\mu$, fix a vertex $v_\mu$ with $\chi(\hub,v_\mu)=\mu$. No vertex,
color, or color neighborhood is supplied as an additional input
predicate.

Jipsen, Maddux, and Tuza~\citep{jipsen-maddux-tuza1995} gave
the random and affine-plane colorings underlying our reference
structures. Although they formulated a one-witness condition, their
constructions also yield the two-witness property used here. The
random-coloring argument above establishes it at every sufficiently
large domain size; their affine-plane coloring has it for planes of
order at least four. Appendix~\ref{app:explicit} verifies the latter
and gives an explicit construction on fewer than $16N^2$ vertices.

\subsection{Color-determined relations, folds, and predictions}\label{sec:types}
For $\mu\in\Lambda$, let
$e_\mu\coloneqq\{(u,v)\in V^2:u\ne v,\ \chi(u,v)=\mu\}$ be the set of edges of
color $\mu$, in both orientations. We extend $\chi$ to loops by giving
every loop the extra label $\diagc$, which is not a color. A relation
$X\subseteq V^2$ is \defn{color-determined} if $X(u,v)$ depends only
on $\chi(u,v)$. For such a relation, write $X[\mu]\in\bool$ for its
common value at pairs $(u,v)$ with $\chi(u,v)=\mu$, where
$\mu\in\Lambda\cup\{\diagc\}$. These $N+1$ values form a table, and
\[
 S_X\coloneqq\{\mu\in\Lambda:X[\mu]=1\}
\]
is the \defn{off-diagonal support} of $X$. Thus $X$ is the union of the
classes $e_\mu$ with $\mu\in S_X$, together with the diagonal if
$X[\diagc]=1$. A reference structure interprets every input relation
by a color-determined relation; we write $\mathcal M_0$ for it.

\begin{lemma}[Closure of color-determined relations]\label{lem:tables}
The constants $\zero,\one,\id$ are color-determined, and
color-determined relations are closed under Boolean operations, converse,
and composition. If $X$ and $Y$ are color-determined, then for every color
$\mu\in\Lambda$,
\begin{equation*}\tag{CT}\label{eq:table-off}
\begin{aligned}
 (X;Y)[\mu]
  &=\bigl(X[\diagc]\land Y[\mu]\bigr)\lor\bigl(X[\mu]\land Y[\diagc]\bigr)
       \lor\bigl(S_X\ne\varnothing\land S_Y\ne\varnothing\bigr),\\
 (X;Y)[\diagc]
  &=\bigl(X[\diagc]\land Y[\diagc]\bigr)\lor\bigl(S_X\cap S_Y\ne\varnothing\bigr),
\end{aligned}
\end{equation*}
where a condition stands for its truth value.
\end{lemma}
\begin{proof}
Boolean operations act entry by entry, and converse preserves this
property because $\chi$ is symmetric. For composition, consider a pair
$(u,v)$ and a potential witness $z$. An endpoint witness,
$z\in\{u,v\}$, contributes the disjuncts of~\eqref{eq:table-off} that involve
$X[\diagc]$ or $Y[\diagc]$. An outside witness, $z\notin\{u,v\}$,
contributes $X[\chi(u,z)]\land Y[\chi(z,v)]$. If $u\ne v$, then by
Lemma~\ref{lem:palette} the color pairs $(\chi(u,z),\chi(z,v))$ of outside
witnesses range over all of $\Lambda\times\Lambda$, so an outside witness
exists exactly when $S_X$ and $S_Y$ are both nonempty. If $u=v$, both
edges of an outside witness are the edge $\{u,z\}$, and every color
occurs on an edge at $u$, so an outside witness exists exactly when
$S_X\cap S_Y\ne\varnothing$.
\end{proof}
In particular, every value of a relational circuit on $\mathcal M_0$ is
color-determined, and
\[
 e_\mu;e_\nu=
 \begin{cases}\one,&\mu=\nu,\\
 \overline{\id},&\mu\ne\nu.
 \end{cases}
\]
Thus the color-determined relations give a concrete representation of
the symmetric relation algebra $E_{N+1}(1,2,3)$ studied by Jipsen,
Maddux, and Tuza~\citep{jipsen-maddux-tuza1995}. Its atoms are $\id$
and the $e_\mu$, and $\overline{\id}=\bigcup_{\mu\in\Lambda}e_\mu$ is
the diversity relation. This identification follows from the displayed
multiplication table, not from automorphisms of the reference graph.

\paragraph{Folds and predictions.}
From now on, take the palette $\Lambda\coloneqq2^{[k]}$ for some $k\ge1$, so that
$N=2^k$. The complement of a color $\mu$ is the color
$\mu^c=[k]\setminus\mu$, and no color equals its complement. The fold of
a color $\lambda$ recolors every edge of color $\lambda$ at the hub to
$\lambda^c$ and leaves all other edges unchanged: for $u\ne v$,
\[
 \chi_\lambda(u,v)\coloneqq
 \begin{cases}\lambda^c,&\hub\in\{u,v\}\text{ and }\chi(u,v)=\lambda,\\
 \chi(u,v),&\text{otherwise},
 \end{cases}
\]
and every loop keeps the label $\diagc$. Only two properties of the fold
matter below: it removes exactly the color $\lambda$ from the hub, and
it changes no edge away from the hub. Any target color other than
$\lambda$ would serve as well. For a color-determined relation $X$,
its prediction $\pred X\lambda$ reads the same values on $\chi_\lambda$:
\[
 \pred X\lambda(u,v)\coloneqq X[\chi_\lambda(u,v)]\qquad(u,v\in V).
\]
Thus $\pred X\lambda$ agrees with $X$ except on the recolored edges,
where it takes the value $X[\lambda^c]$. Replacing every input relation
by its prediction defines the folded structure $\mathcal M_\lambda$ on
the same domain $V$. Predictions are defined only for color-determined relations, in
particular for all values computed on $\mathcal M_0$. We do not assume
that relations computed on $\mathcal M_\lambda$ are color-determined;
in general, they are not.

\subsection{Noticing, and unnoticed folds}\label{sec:unnoticed}
Let $\lambda$ be a color and $\mathbf U=(U_1,\ldots,U_r)$ a tuple of
color-determined relations. An $r$-ary operation $\operatorname{op}$
notices the fold of $\lambda$ at $\mathbf U$ if
\[
 \operatorname{op}\bigl(\pred{U_1}\lambda,\ldots,\pred{U_r}\lambda\bigr)\ne
 \pred{\operatorname{op}(\mathbf U)}\lambda.
\]
Both sides are determined by $\mathbf U$ and $\lambda$: noticing is a
property of the operands on $\mathcal M_0$, not of any computation on
$\mathcal M_\lambda$.

\begin{proposition}[Commutation with prediction]\label{prop:prediction-commutation}
For any color $\lambda$ and color-determined relations $X,Y$, prediction
commutes with Boolean operations and converse:
\[
\begin{aligned}
 \pred{(X\cup Y)}\lambda&=\pred X\lambda\cup\pred Y\lambda,&
 \pred{(X\cap Y)}\lambda&=\pred X\lambda\cap\pred Y\lambda,\\
 \pred{(\overline X)}\lambda&=\overline{\pred X\lambda},&
 \pred{(X^\smile)}\lambda&=(\pred X\lambda)^\smile.
\end{aligned}
\]
The constants $\zero,\one,\id$ are also fixed by prediction.
\end{proposition}
\begin{proof}
Boolean operations act pair by pair, so both sides apply the same
operation to the entries at $\chi_\lambda(u,v)$. Converse commutes with
prediction because $\chi_\lambda(u,v)=\chi_\lambda(v,u)$. The constants
are fixed because a fold preserves the diagonal.
\end{proof}

Thus only composition can notice a fold. A composition gate of a
relational circuit notices the fold of $\lambda$ if the composition of
its operand values on $\mathcal M_0$ notices it.
\begin{lemma}[Unnoticed folds]\label{lem:unnoticed}
Let $C$ be a relational circuit and $\lambda$ a color. If no composition
gate of $C$ notices the fold of $\lambda$, then on $\mathcal M_\lambda$
every gate of $C$ computes the prediction of its value on $\mathcal M_0$.
In particular, the outputs of $C$ on $\mathcal M_0$ and on
$\mathcal M_\lambda$ agree at every pair that the fold does not recolor.
\end{lemma}
\begin{proof}
Evaluate $C$ on $\mathcal M_0$. By Lemma~\ref{lem:tables}, every gate
value is color-determined, so its prediction is defined. We show by
induction along a topological order of the gates that every gate computes
on $\mathcal M_\lambda$ the prediction of its value on $\mathcal M_0$. For
input relations this is the definition of $\mathcal M_\lambda$, and
constants are their own predictions by Proposition~\ref{prop:prediction-commutation}.
For Boolean and converse gates, the claim follows from the induction
hypothesis and Proposition~\ref{prop:prediction-commutation}.
If the operands of a composition gate have values $X$ and $Y$ on
$\mathcal M_0$, then by induction they have values $\pred X\lambda$ and $\pred Y\lambda$ on
$\mathcal M_\lambda$. The gate therefore computes
$\pred X\lambda;\pred Y\lambda$ on $\mathcal M_\lambda$, which equals
$\pred{(X;Y)}\lambda$ because the gate does not notice the fold. Finally,
every value $U$ on $\mathcal M_0$ agrees with $\pred U\lambda$ at every
pair that the fold does not recolor.
\end{proof}

\subsection{Each composition notices at most two folds}\label{sec:responses}
Fix color-determined relations $X,Y$ and a color $\lambda$, and write
$J\coloneqq\Lambda\setminus\{\lambda\}$ for the set of colors at $\hub$
after the fold of $\lambda$.
Endpoint witnesses contribute the same value to
$\pred X\lambda;\pred Y\lambda$ and to $\pred{(X;Y)}\lambda$. At a pair
$(u,v)$ with $u\ne v$ and $\mu\coloneqq\chi_\lambda(u,v)$, they contribute
$(X[\diagc]\land Y[\mu])\lor(X[\mu]\land Y[\diagc])$ to
$\pred X\lambda;\pred Y\lambda$, because $\pred X\lambda$ and
$\pred Y\lambda$ take the values $X[\diagc]$ and $Y[\diagc]$ on the
diagonal and $X[\mu]$ and $Y[\mu]$ at $(u,v)$, and the same value to
$\pred{(X;Y)}\lambda(u,v)=(X;Y)[\mu]$ by~\eqref{eq:table-off}; at a
diagonal pair, they contribute $X[\diagc]\land Y[\diagc]$ to both. The
following table compares the outside witnesses. Its left column states
when $\pred X\lambda;\pred Y\lambda$ has an outside witness at the pair,
and its right column states the last disjunct of~\eqref{eq:table-off},
which $\pred{(X;Y)}\lambda$ reads at the pair; here $u$ and $v$ range
over vertices other than $\hub$, and $S\ne\varnothing\ne S'$ means that
both sets are nonempty.
\begin{center}
\renewcommand{\arraystretch}{1.15}
\begin{tabular}{lll}
\toprule
Output pair & Outside witness in $\pred X\lambda;\pred Y\lambda$
 & Last disjunct of~\eqref{eq:table-off}\\
\midrule
$(u,v)$, $u\ne v$ & $S_X\ne\varnothing\ne S_Y$ & $S_X\ne\varnothing\ne S_Y$\\
$(u,u)$ & $S_X\cap S_Y\ne\varnothing$ & $S_X\cap S_Y\ne\varnothing$\\
$(\hub,v)$ & $S_X\cap J\ne\varnothing\ne S_Y$ & $S_X\ne\varnothing\ne S_Y$\\
$(u,\hub)$ & $S_X\ne\varnothing\ne S_Y\cap J$ & $S_X\ne\varnothing\ne S_Y$\\
$(\hub,\hub)$ & $S_X\cap S_Y\cap J\ne\varnothing$ & $S_X\cap S_Y\ne\varnothing$\\
\bottomrule
\end{tabular}
\end{center}
For the left column, consider the outside witnesses $z$ of a pair. For
distinct $u,v\ne\hub$, Lemma~\ref{lem:palette} gives, for every pair of
colors, an outside witness other than $\hub$, whose two edges keep their
colors; the witness $\hub$ contributes one more color pair, which is
already realized. For $(u,u)$ with $u\ne\hub$, every color occurs on an
edge at $u$ that avoids $\hub$. At a pair $(\hub,v)$ of the
\defn{hub row}, the left edge $\{\hub,z\}$ of every outside witness lies
at the hub, so its color becomes $\chi_\lambda(\hub,z)\in J$, while the
right edge $\{z,v\}$ keeps its color; by Lemma~\ref{lem:palette} the
realized color pairs are exactly $J\times\Lambda$. The pairs $(u,\hub)$
of the \defn{hub column} are symmetric, with $\Lambda\times J$. At
$(\hub,\hub)$ both edges of an outside witness are the edge
$\{\hub,z\}$, whose color $\chi_\lambda(\hub,z)$ ranges exactly over $J$.

\begin{lemma}[When a composition notices a fold]\label{lem:local}\label{lem:notice}
Let $X,Y$ be color-determined. The composition $X;Y$
notices the fold of $\lambda$ only if
\[
 S_X=\{\lambda\},\qquad S_Y=\{\lambda\},\qquad\text{or}\qquad
 S_X\cap S_Y=\{\lambda\}.
\]
Hence it notices at most two folds, at most one if $S_X$ or $S_Y$ is
$\varnothing$ or $\Lambda$, and none if both are.
\end{lemma}
\begin{proof}
By the table, the two sides can differ only in the hub row, in the hub
column, or at $(\hub,\hub)$. In the hub row they can differ only if
$S_X\ne\varnothing$ but $S_X\cap J=\varnothing$, that is,
$S_X=\{\lambda\}$; in the hub column only if $S_Y=\{\lambda\}$; and at
$(\hub,\hub)$ only if $S_X\cap S_Y=\{\lambda\}$. Each condition
determines $\lambda$. If $S_X\cap S_Y\ne\varnothing$, each of the first
two conditions implies the third, so at most one fold is noticed;
otherwise the third condition never holds, and at most two are. If $S_Y$
is $\varnothing$ or $\Lambda$, then $S_Y\ne\{\lambda\}$ because $N\ge2$,
and $S_X\cap S_Y$ is $\varnothing$ or $S_X$, so only the condition
$S_X=\{\lambda\}$ remains; the case of $S_X$ is symmetric. If both are
$\varnothing$ or $\Lambda$, none of the conditions can hold.
\end{proof}
The bound two is attained. The composition $e_\mu;e_{\mu^c}$ has
$S_X=\{\mu\}$ and $S_Y=\{\mu^c\}$, and notices the folds of both $\mu$
and $\mu^c$. On the reference structure of \S\ref{sec:G}, the relation
$e_\mu$ is the relation $H_\mu$ of the DNF.
Proposition~\ref{prop:promise} computes $G_k$ on the reference structure
and all its folded structures with $2^{k-1}$ such compositions.
By Lemma~\ref{lem:notice}, each composition gate notices at most two
folds. What a gate notices depends only on its operand values on
$\mathcal M_0$, so a shared gate counts once, regardless of its fan-out.
Hence at most $2\gcost C$ folds are noticed by some gate of $C$, and
Lemma~\ref{lem:unnoticed} applies to all other folds. Nothing here
depends on the query or on the interpretation of the input relations. If
a query has different answers on $\mathcal M_0$ and on
$\mathcal M_\lambda$ at a pair that the fold of $\lambda$ does not
recolor, for every color $\lambda$, then every relational circuit $C$
defining it satisfies $2\gcost C\ge2^k$. Section~\ref{sec:logical} shows
that $G_k$ is such a query.

\subsection{Tests and the test--cost theorem}\label{sec:simultaneous}
For randomized circuits and for the other targets of
\S\ref{sec:additional-results}, we use a counting form of this argument
that refers only to a few abstract properties. For a color $\lambda$,
the \defn{test} for $\lambda$ consists of the structures $\mathcal M_0$
and $\mathcal M_\lambda$ and the \defn{output observation}
$\rho_\lambda(U)\coloneqq U(\hub,v_{\lambda^c})$, the entry of an output
value $U$ at the pair $(\hub,v_{\lambda^c})$ of \S\ref{sec:structure}.
Since $\lambda^c\ne\lambda$, the fold of $\lambda$ does not recolor this
pair, so
\begin{equation*}\tag{O}\label{eq:observation}
 \rho_\lambda(\pred U\lambda)=\rho_\lambda(U)
\end{equation*}
for every reference output value $U$, as defined in~(i) below.

The counting argument uses only three properties of this construction.
We call data satisfying them a \defn{test system}:
\begin{enumerate}[label=(\roman*),leftmargin=*]
\item Each type of value has a collection of \defn{reference values},
which contains the input values on $\mathcal M_0$ and is closed under the
permitted operations, and every reference value $U$ has a prediction
$\pred U\lambda$ for each color $\lambda$.
\item Every input value on $\mathcal M_\lambda$ is the prediction of the
same input value on $\mathcal M_0$.
\item Every test satisfies~\eqref{eq:observation}, and the required
answers at its observation on $\mathcal M_0$ and on $\mathcal M_\lambda$
differ.
\end{enumerate}
For relational circuits, the reference values are the color-determined
relations, and Lemma~\ref{lem:tables} supplies~(i). The typed scalar
computations in \S\ref{sec:matrix} will satisfy the same conditions;
predictions carry the type of their argument, which we suppress in the
notation. For a gate $g$ of a deterministic circuit $C$, let
$\operatorname{op}_g$ be its operation and $\mathbf U_g$ the tuple of its
operand values on $\mathcal M_0$. As in \S\ref{sec:unnoticed}, $g$
notices the fold of $\lambda$ if
$\operatorname{op}_g(\pred{\mathbf U_g}\lambda)\ne
\pred{\operatorname{op}_g(\mathbf U_g)}\lambda$, where predictions act
componentwise on tuples. Define
\[
 \begin{split}
 D_g&\coloneqq\{\lambda\in\Lambda:g\text{ notices the fold of }\lambda\},\\
 \tcost(C)&\coloneqq\sum_{g\text{ distinct}}|D_g|.
 \end{split}
\]
For input nodes put $D_g\coloneqq\varnothing$; constant gates are nullary
operations. The definition of $D_g$ compares the entire output of the
gate, and it uses predicted operands computed from $\mathcal M_0$, not
operands recomputed on $\mathcal M_\lambda$.

We call $\tcost(C)$ the \defn{cost} of $C$. It counts the pairs of a
distinct gate and a fold that the gate notices. This integer depends on
the reference structure and the tests, which remain fixed, but is
defined even for an incorrect circuit.

\begin{theorem}[Test--cost theorem]\label{thm:test-cost}
Fix a test system. A deterministic circuit $C$ correct at both
observations of every test satisfies
\[
 2^k\le\tcost(C).
\]
If $C_\omega$ is chosen from an input-independent distribution of circuits,
and each observation in each test is incorrect with probability at most
$\varepsilon<1/2$, then
\[
 (1-2\varepsilon)2^k\le\mathbb E_\omega\tcost(C_\omega).
\]
\end{theorem}
\begin{proof}
Fix a deterministic circuit, not assumed correct, and put
$D_C\coloneqq\bigcup_gD_g$. For $\lambda\notin D_C$, no gate notices the
fold of $\lambda$, so the induction in the proof of
Lemma~\ref{lem:unnoticed}, with~(ii) for the input values, gives
\[
 \sem C{\mathcal M_\lambda}=\pred{(\sem C{\mathcal M_0})}\lambda.
\]
The equality concerns the whole output. By the observation
condition~\eqref{eq:observation}, the two observations are equal, whereas
their correct values differ. Hence at least one is wrong. If
$\mathrm{err}_{\lambda,0}(C),\mathrm{err}_{\lambda,1}(C)$ are the
corresponding zero-one error indicators, then
\[
 \tcost(C)+\sum_{\lambda\in\Lambda}
   \bigl(\mathrm{err}_{\lambda,0}(C)+\mathrm{err}_{\lambda,1}(C)\bigr)\ge2^k,
\]
because $|D_C|\le\tcost(C)$. Zero error proves the first assertion.
Take expectations using the same random circuit on the two structures
of each test. Each error indicator has expectation at most $\varepsilon$,
which proves the second assertion. No independence among the tests or
their errors is required.
\end{proof}

This is the counting argument for all subsequent lower bounds on the
reference structure. Each of them bounds $\tcost(C)$ from above by a
resource count of the target, such as $2\gcost C$ in
Lemma~\ref{lem:cor-cost} below, and therefore also holds in two further
forms, which we do not restate. First, Lemma~\ref{lem:palette} gives the
reference structure on every domain of size $n\ge n_0(2^k)$, and the
folds keep the domain, so each bound holds for targets that need to be
correct only on $n$-element inputs and may be chosen for $n$. Second, by
the randomized assertion of Theorem~\ref{thm:test-cost}, each bound for
targets that compute a query also holds for an input-independent
distribution of targets with pointwise error at most
$\varepsilon<1/2$, after the resource count is replaced by its
expectation and the bound is multiplied by $1-2\varepsilon$; for scalar
outputs, an error is a failure to return the exact required value.

For ordinary \CoR, call a composition gate whose operand values on
$\mathcal M_0$ are $X,Y$ \defn{one-sided} if $S_X$ or $S_Y$ is
$\varnothing$ or $\Lambda$. Let
$m_1(C)$ count these gates and $m_2(C)$ the other composition gates,
counting shared gates once.
\begin{lemma}[Cost of compositions]\label{lem:cor-cost}
Every relational circuit $C$ satisfies
\[
 \tcost(C)\le m_1(C)+2m_2(C)\le2\gcost{C}.
\]
\end{lemma}
\begin{proof}
By Lemma~\ref{lem:local}, a one-sided composition notices at most one
fold and any other composition at most two. Other gates notice nothing.
Summing over distinct gates proves the inequalities.
\end{proof}
The projections $X;\one$ and $\one;X$ are one-sided. If a syntactic count
is preferred, classify only these projections as one-sided and put every
remaining composition in $m_2$; Lemma~\ref{lem:cor-cost} still holds.
